\chapter*{Preface}
\addcontentsline{toc}{chapter}{Preface}\markboth{\MakeUppercase{Preface}}{\MakeUppercase{Preface}}

The Riemann Hypothesis can be stated as the positivity of a measure. The reciprocal Xi clock $u\mapsto\xi(\frac12)/\xi(\frac12+\sqrt u)$ is the Laplace transform of a generalised gamma convolution exactly when every nontrivial zero lies on the critical line, and then its Thorin measure is the discrete measure $\sum_{\gamma>0}\delta_{\gamma^2}$ on the squared ordinates. This book asks what computation can say about that positivity condition.

Its answer has three parts. First, the Thorin measure is computable at every positive resolution, with explicit error bounds, from the P\'olya kernel alone, with no primes and no list of zeros. Second, the distance between the computed measure and the measure demanded by RH can be measured exactly, and the signature an off line zero would leave can be calibrated. Third, a finite computation could refute RH by certifying a zero off the line, but neither a computation of the zeros up to a finite height nor a finite set of positivity tests on the heat trace can prove it: for every sufficiently large height $R$ such that RH has been verified up to $R+50$, there is a perturbation of $\xi$, with the same order, type and symmetry and the same zeros below $R$, that has zeros off the line between $R$ and $R+50$ and still passes the tests. The book makes this precise and locates the remaining difficulty. A last question is how much of the positivity can be proved. Moved from the centre of the critical strip to the basepoint one, the clock has a positive Thorin measure with no hypothesis. Chapters~\ref{chap:basepoint} to \ref{chap:almost} follow each construction of the book back across the strip as far as it holds unconditionally; below the basepoint one, zero density estimates still show that the zeros to the right of the basepoint, the only possible source of negativity, carry a vanishing fraction of the mass of the density.

The book grew out of the first author's programme on generalised gamma convolutions and the P\'olya density, collected in the monograph \emph{On Hilbert's 8th Problem} \cite{P26}, and out of a joint paper computing the Thorin measure directly \cite{PS26}. The first author's arXiv paper of the same title is a separate work, which we cite by version: \cite{P17} is its version 24, of 2026, and \cite{P17v5} its version 5, of 2017. Chapters~\ref{chap:rzeta} and \ref{chap:rzid} incorporate Part IX of that monograph \cite{P26}, the first author's work on the reciprocal zeta function, written as the manuscript \cite{P26rz}, and Chapters~\ref{chap:basepoint} to \ref{chap:almost} a paper of the authors on the reciprocal Xi function at the basepoint one and across the critical strip, together with notes on Riesz's criterion and on almost generalised gamma convolutions. Sections~\ref{ch:chi12} and \ref{ch:bpeta} draw on a manuscript of the first author on the Thorin measures of the Dedekind eta function \cite{P26eta}, and Chapter~\ref{chap:ferro} on a note of the first author on ferromagnetic random variables. The core evaluator runs in ball arithmetic, and Appendix~\ref{app:code} describes the scripts and data files behind the tables and figures and what the archive lacks. No proof of RH is claimed.

\bigskip
\noindent Chicago and Fairfax, September 2026
